\begin{frame}{}
  \begin{block}{Teorema}
    Seja
    \begin{align*}
      x^{\prime}=f(x),\quad x\in U\subseteq\mathbb{R}^{n},
    \end{align*}
    um sistema autónomo com $f\in C^{1}(U,\mathbb{R}^{n})$.\\
    Seja $p\in U$ singularidade hiperbólica do sistema. Então, existem vizinhanças $V$ de $p$, $W$ de $0\in\mathbb{R}^{n}$ e um homeomorfismo $h:V\to W$ tal que
    \begin{enumerate}
      \item $h(p)=0$.
      \item Se $\phi$ é o sistema dinámico de $x^{\prime}=f(x)$ e considerando $y^{\prime}=Ay$ com $A=Df(p)$, então
        \begin{align*}
          h(\phi_t(x))=e^{tA}h(x),\quad\text{para todo }\phi_t(x)\in V.
        \end{align*}
    \end{enumerate}
    Isto é, $h$ é uma conjugação topológica local dos sitemas $x^{\prime}=f(x)$ e $y^{\prime}=Ay$. 
  \end{block} 
\end{frame}

\begin{frame}{Demonstração}
  \begin{block}{}
    Considere $T(x)=x+p$. Seja $\tilde{U}=T^{-1}(U)$ e $\tilde{f}:\tilde{U}\to\mathbb{R}^{n}$ dada por $\tilde{f}(x)=f(T(x))$, logo temos que
    \begin{enumerate}
      \item $\tilde{f}(0)=f(p)=0$ e $D\tilde{f}(0)=Df(p)$.
      \item $\tilde{f}\in C^{1}(\tilde{U},\mathbb{R}^{n})$.
      \item $0$ é singularidade hiperbólica de $x^{\prime}=\tilde{f}$. 
    \end{enumerate}
    Portanto, sem perdida de generalidade, vamos supor  que $p=0$.
  \end{block}
  \begin{block}{}
    Definindo $A=Df(0)$ e $r(x)=f(x)-Ax$, temos que
    \begin{enumerate}
      \item $x^{\prime}=f(x)=Ax+r(x)$.
      \item $r(0)=f(0)-A(0)=0$.
      \item $r\in C^{1}(U,\mathbb{R})$.
      \item $Dr(x)=Df(x)-A$, logo $Dr(0)=Df(0)-A=0$.
    \end{enumerate}
  \end{block}
\end{frame}

\begin{frame}{}
  \begin{block}{}
    Daí temos que
    \begin{enumerate}
      \item $\lim_{x \to 0}\frac{\norm{r(x)}}{\norm{x}}=0$.
      \item $\lim_{x \to 0}r(x)=0$.
      \item $\lim_{x \to 0}Dr(x)=0$.
    \end{enumerate}
    Assim, dado $\epsilon>0$ temos que existe $\delta>0$ tal que se $\norm{x}<\delta$, então $\norm{Dr(x)}<\epsilon$. Logo pela Desigualdade do Valor Meio temos que se $x,y\in B_{\delta}(0)$, então
    \begin{equation}\label{eq:r-lip}
      \norm{r(x)-r(y)}<\epsilon\norm{x-y}
    \end{equation}
    Pelo que temos que $r\big|_{B_{\delta}(0)}$ é Lipschitz. 
  \end{block}
\end{frame}

\begin{frame}{}
  \begin{block}{}
    Agora, considere $\eta^{\infty}_c(\mathbb{R}^{n},[0,1])$ tal que
    \begin{align*}
      \eta(x)=\begin{cases}
        1, &\text{ se } |x|<\frac{1}{2} \text{,} \\
        0, &\text{ se } |x|\geq 1.
      \end{cases}
    \end{align*}
    Definindo $C=\sup_{x\in\mathbb{R}^{n}}\eta^{\prime}(x)<+\infty$. Note que se definimos $\chi_{\delta}(x)=\eta\left( \frac{x}{\delta} \right)$ temos que
    \begin{equation}\label{eq:cota-xlinha}
      |\chi_{\delta}^{\prime}(x)|\leq \frac{1}{\delta}\left| \eta\left( \frac{x}{\delta} \right) \right|\leq\frac{c}{\delta}.
    \end{equation}
    Por outro lado temos que
    \begin{equation}\label{eq:cota-rx}
      \norm{r(x)}=\norm{r(x)-r(0)}< \epsilon\norm{x}\leq \epsilon\delta.
    \end{equation}
  \end{block} 
\end{frame}

\begin{frame}{}
  \begin{block}{}
    Sendo assim, defina $\tilde{r}(x)=\chi_{\delta}(\norm{x})\cdot r(x)$, logo
    \begin{align*}
      \tilde{r}(x)=\begin{cases}
        r(x), &\text{ se } \norm{x}\leq \frac{\delta}{2} \text{,} \\
        0, &\text{ se } \norm{x}>\delta.
      \end{cases}
    \end{align*}
    Além disso, fazendo ums cálculos e usando $\eqref{eq:r-lip}$, \eqref{eq:cota-xlinha} e $\eqref{eq:cota-rx}$ temos que
    \begin{align*}
      \norm{D\tilde{r}(x)}\leq \frac{c}{\delta}\epsilon\delta+\epsilon=(c+1)\epsilon=\tilde{\epsilon}.
    \end{align*}
    Pelo que temos que $\tilde{r}$ é Lipschitz com $Lip(\tilde{r})<\tilde{\epsilon}$. 
  \end{block}
\end{frame}
